\documentclass[11pt]{article}
\usepackage[margin=1in]{geometry}
\usepackage{amsmath,amssymb,amsthm}
\usepackage[T1]{fontenc}
\usepackage{lmodern}
\usepackage{graphicx}
\usepackage[hidelinks]{hyperref}

\newtheorem{theorem}{Theorem}
\newtheorem{lemma}{Lemma}
\newtheorem{proposition}{Proposition}

\title{Candidate Full Result: The Random-Semigroup SLLN Holds on Every Banach Space}
\author{Run \texttt{fa\_banach\_001}, agent \texttt{agent\_lane\_11}}
\date{2026-06-26}

\newcommand{\X}{\mathcal X}
\newcommand{\Lcal}{\mathcal L}
\newcommand{\EE}{\mathbb E}
\newcommand{\PP}{\mathbb P}
\newcommand{\N}{\mathbb N}
\newcommand{\norm}[1]{\left\lVert #1\right\rVert}
\newcommand{\barA}{\overline A}
\emergencystretch=2em

\begin{document}
\maketitle

\paragraph{Status.}
\textbf{Candidate full solution, likely valid, needs human review.}
This packet proves Conjecture 2.6 of Dzhenzher--Sakbaev
\cite{DzhenzherSakbaev2025}: the strong-operator strong law for bounded
i.i.d. random semigroup generators does not require uniform smoothness of the
Banach space.

\section*{Source Conjecture}

The source paper proves the result in strong operator topology when \(\X\) is
uniformly smooth, and in weak operator topology for arbitrary Banach spaces.
On page 4, after noting the \(\ell_1\) consequence, it states Conjecture 2.6:

\begin{center}
\includegraphics[width=\textwidth]{figures/open_problem_crop.png}
\end{center}

Transcribed, the conjecture says:

\begin{quote}
Let \(\X\) be a Banach space. Then under the assumptions of Theorem 2.4,
\(e^{A_1t/n}\cdots e^{A_nt/n}\) converges a.s. in SOT of
\(\norm{\cdot}_{\X}\) to \(e^{\EE A t}\) uniformly for \(t>0\) in any segment.
\end{quote}

Here \(A,A_1,A_2,\ldots\) are i.i.d. bounded random elements of
\(\Lcal(\X)\), and \(\EE A\) is the Bochner expectation.

\section*{Proof Intuition}

The source proof uses uniform smoothness to obtain summable martingale moment
estimates for the random product.  For bounded generators there is a more
direct pathwise route.  First replace each exponential
\(e^{A_jt/n}\) by the Euler factor \(I+(t/n)A_j\); the accumulated replacement
error is \(O(1/n)\), uniformly for \(t\) in a compact interval.  The Euler
product is then compared with the deterministic Euler product for
\(\bar A=\EE A\).

For a fixed \(t\), the accumulated error solves a linear recursion driven by
mean-zero increments
\[
  \frac{t}{n}(A_j-\bar A)y_{n,j},
\]
where \(y_{n,j}\) lies on the deterministic Euler orbit of \(x\).  That orbit
is contained in a compact curve.  A finite partition of the curve reduces the
noise estimate to finitely many ordinary Banach-valued strong laws of large
numbers, applied on intervals of indices.  A discrete summation-by-parts
estimate then converts small partial sums of the noise into small product
error.  Finally, a uniform Lipschitz bound in \(t\) upgrades convergence on
rational \(t\)'s to uniform convergence on every compact interval.

\section*{Statement}

\begin{theorem}
\label{thm:all-banach-slln}
Let \(\X\) be a Banach space. Let
\(A,A_1,A_2,\ldots\) be i.i.d. Bochner-measurable random elements of
\(\Lcal(\X)\), and assume that \(\norm{A}\le \rho\) almost surely for some
\(\rho<\infty\).  Put \(\bar A=\EE A\).  Then for every \(x\in \X\) and every
\(T>0\),
\[
  \lim_{n\to\infty}
  \sup_{0\le t\le T}
  \norm{\left(e^{A_1t/n}\cdots e^{A_nt/n}-e^{\bar A t}\right)x}
  =0
\]
almost surely.
\end{theorem}

\section*{Proof}

We use the following elementary form of the Banach-valued strong law.  It is a
standard consequence of the usual SLLN for Bochner-integrable i.i.d.
Banach-valued random variables; boundedness is more than enough for the
application below.

\begin{lemma}[Interval SLLN]
\label{lem:interval-slln}
Let \(Z_1,Z_2,\ldots\) be i.i.d. bounded random variables with values in a
Banach space \(E\), and assume \(\EE Z_1=0\).  Then, almost surely,
\[
  \sup_{0\le a\le b\le n}
  \norm{\frac1n\sum_{j=a+1}^{b} Z_j}_E \longrightarrow 0 .
\]
\end{lemma}

\begin{proof}
The ordinary Banach-valued SLLN gives
\(m^{-1}\sum_{j=1}^m Z_j\to0\) almost surely.  Hence
\[
  \max_{0\le m\le n}
  \norm{\frac1n\sum_{j=1}^{m}Z_j}_E\to0
\]
almost surely: finitely many initial partial sums are divided by \(n\), while
the remaining ones are controlled by the SLLN.  The interval estimate follows
by subtracting two partial sums.
\end{proof}

\begin{lemma}[Compact-curve noise estimate]
\label{lem:curve-slln}
Let \(B_j=A_j-\bar A\).  Fix \(x\in\X\) and \(t\ge0\).  Let \(h_n=t/n\) and
\[
  y_{n,k}=(I+h_n\bar A)^k x,\qquad 0\le k\le n .
\]
Then, almost surely,
\[
  \max_{0\le m\le n}
  \norm{\frac1n\sum_{k=1}^{m} B_{n-k+1}y_{n,k-1}}
  \longrightarrow0 .
\]
\end{lemma}

\begin{proof}
The points \(y_{n,k}\) converge uniformly, as sampled Euler points, to the
compact curve \(s\mapsto e^{s\bar A}x\), \(0\le s\le t\).  Given
\(\varepsilon>0\), choose a partition
\[
  0=s_0<s_1<\cdots<s_M=t
\]
so fine that both \(e^{s\bar A}x\) and the Euler samples vary by at most
\(\varepsilon\) inside each subinterval, for all sufficiently large \(n\).
For \(k\) with \((k-1)t/n\in[s_{r-1},s_r)\), replace \(y_{n,k-1}\) by the
fixed vector \(z_r=e^{s_{r-1}\bar A}x\).  Since
\(\norm{B_j}\le2\rho\), the total error in the normalized sums is at most
\(2\rho\varepsilon\).

It remains to handle finitely many fixed vectors \(z_r\).  For each \(r\), the
variables \(B_jz_r\) are i.i.d., bounded, and have mean zero.  Lemma
\ref{lem:interval-slln} gives
\[
  \sup_{0\le a\le b\le n}
  \norm{\frac1n\sum_{j=a+1}^{b} B_jz_r}\to0
\]
almost surely.  The sums appearing after the replacement are finite sums of
such interval sums, because \(k\mapsto n-k+1\) reverses intervals.  Intersect
the finitely many probability-one events and let \(\varepsilon\downarrow0\).
\end{proof}

\begin{lemma}[Euler products]
\label{lem:euler}
For each fixed \(t\ge0\),
\[
  \left(I+\frac{t}{n}A_1\right)\cdots
  \left(I+\frac{t}{n}A_n\right)x
  -
  \left(I+\frac{t}{n}\bar A\right)^n x
  \longrightarrow0
\]
almost surely.
\end{lemma}

\begin{proof}
Write \(h=t/n\), \(C_j=I+hA_j\), and \(\bar C=I+h\bar A\).  To match the
source product order, define
\[
  X_{n,0}=x,\qquad
  X_{n,k}=C_{n-k+1}X_{n,k-1},\quad 1\le k\le n,
\]
so \(X_{n,n}=C_1\cdots C_nx\).  Also put \(y_{n,k}=\bar C^kx\) and
\(D_{n,k}=X_{n,k}-y_{n,k}\).  Then
\[
  D_{n,k}=C_{n-k+1}D_{n,k-1}
  +h(A_{n-k+1}-\bar A)y_{n,k-1}.
\]
Set
\[
  \eta_{n,k}=h(A_{n-k+1}-\bar A)y_{n,k-1},
  \qquad
  S_{n,m}=\sum_{k=1}^{m}\eta_{n,k}.
\]
By Lemma \ref{lem:curve-slln},
\(\max_{m\le n}\norm{S_{n,m}}\to0\) almost surely.

Solving the linear recursion gives
\[
  D_{n,m}=\sum_{k=1}^{m}R_{m,k}\eta_{n,k},
\]
where \(R_{m,m}=I\) and
\[
  R_{m,k}=C_{n-m+1}\cdots C_{n-k}
  \qquad (k<m).
\]
The exact indexing is unimportant; only the following bounds matter:
\[
  \norm{R_{m,k}}\le (1+\rho h)^{m-k}\le e^{\rho t}
\]
and
\[
  \norm{R_{m,k}-R_{m,k+1}}
  \le \rho h e^{\rho t}.
\]
By Abel summation,
\[
  D_{n,m}
  =
  R_{m,m}S_{n,m}
  +\sum_{k=1}^{m-1}(R_{m,k}-R_{m,k+1})S_{n,k}.
\]
Consequently
\[
  \max_{m\le n}\norm{D_{n,m}}
  \le
  \left(1+\rho t e^{\rho t}\right)
  \max_{m\le n}\norm{S_{n,m}},
\]
which tends to zero almost surely.  Taking \(m=n\) proves the lemma.
\end{proof}

\begin{proof}[Proof of Theorem \ref{thm:all-banach-slln}]
First fix \(t\in[0,T]\).  Lemma \ref{lem:euler} gives almost sure convergence
of the random Euler product to \((I+(t/n)\bar A)^nx\), and the latter
converges in norm to \(e^{t\bar A}x\).

It remains to pass from Euler factors to exponentials and then make the
convergence uniform in \(t\).  Since \(\norm{A_j}\le\rho\),
\[
  \norm{e^{hA_j}-(I+hA_j)}
  \le \frac{\rho^2h^2}{2}e^{\rho h}.
\]
Telescoping the product and using
\(\norm{e^{hA_j}}\le e^{\rho h}\), \(\norm{I+hA_j}\le 1+\rho h\), gives
\[
  \sup_{0\le t\le T}
  \norm{
    \left(e^{A_1t/n}\cdots e^{A_nt/n}
    -
    \left(I+\frac tn A_1\right)\cdots
    \left(I+\frac tn A_n\right)\right)x}
  \le \frac{C_{\rho,T}\norm{x}}{n}.
\]
The same deterministic estimate gives
\[
  \sup_{0\le t\le T}
  \norm{\left(I+\frac tn\bar A\right)^nx-e^{t\bar A}x}\to0 .
\]
Thus the exponential product converges almost surely for every fixed rational
\(t\in[0,T]\), after intersecting countably many probability-one events.

Finally, the maps
\[
  t\mapsto e^{A_1t/n}\cdots e^{A_nt/n}x
  \quad\text{and}\quad
  t\mapsto e^{t\bar A}x
\]
are Lipschitz on \([0,T]\) with a common deterministic bound
\[
  \rho e^{\rho T}\norm{x}.
\]
Indeed, differentiate the products and sum the \(n\) terms, each contributing
at most \((\rho/n)e^{\rho T}\norm{x}\).  Convergence on a dense rational set
therefore implies uniform convergence on \([0,T]\).  This proves the
conjecture.
\end{proof}

\section*{Verification Notes}

The proof uses only boundedness, Bochner measurability, independence, and
identical distribution of the generators.  Uniform smoothness is not used.
The argument does not settle the different issue mentioned immediately before
Theorem 2.5 in the source paper: whether the i.i.d. assumption can be weakened
to independent uniformly bounded generators with a common expectation.

A small finite-dimensional smoke script in \texttt{code/matrix\_smoke.py}
checks the theorem numerically for random \(2\times2\) matrix generators.  It
is only a sanity check; the proof above is analytic.

\section*{Novelty Check}

The local run indexes were searched for arXiv `2507.07658`, the source title,
`random semigroups`, `strong law of large numbers`, `SOT`, and related
Banach-space keywords.  No prior packet or attempt for this conjecture was
found.  Web searches on 2026-06-26 for the exact title, `Conjecture 2.6 random
semigroups Banach space`, `SLLN in SOT random semigroups Banach space`, and
close variants found the source arXiv paper and a 2026 follow-up by the same
authors on unbounded generators in uniformly smooth Banach spaces, but no
paper resolving the bounded-generator conjecture for arbitrary Banach spaces.

\section*{Human Review Recommendation}

Review should focus on Lemma \ref{lem:curve-slln}, especially the interval
SLLN reduction for the reversed index order, and on the Abel summation step in
Lemma \ref{lem:euler}.  If those two estimates pass, the packet gives a full
positive solution of Conjecture 2.6.

\begin{thebibliography}{9}

\bibitem{DzhenzherSakbaev2025}
S. V. Dzhenzher and V. Zh. Sakbaev,
\emph{The Strong Law of Large Numbers for random semigroups on uniformly
smooth Banach spaces},
arXiv:2507.07658, 2025.

\bibitem{LedouxTalagrand1991}
M. Ledoux and M. Talagrand,
\emph{Probability in Banach Spaces: Isoperimetry and Processes},
Springer, 1991.

\end{thebibliography}

\end{document}
